\begin{lemma}[Extrema averaging contracts variance]
Let $n>0$, let $k\ge0$ be an integer, and let $a_1,\ldots,a_n$ be
nonnegative real numbers. Write $c=(\sum_t a_t)/n$. There exists a
vector $b$ with the following operation certificate: either $a$ is
constant and $b=a$, or there are distinct indices $i,j$ such that
\[
 a_i\le a_t\le a_j\quad\text{for every }t,\qquad a_i<a_j,
 \qquad
 b=\operatorname{update}(\operatorname{update}(a,i,m),j,m),
 \quad m=(a_i+a_j)/2.
\]
Here $\operatorname{update}(a,i,m)$ replaces entry $i$ by $m$ and
leaves every other entry unchanged. Thus in the nonconstant case $b$
is exactly the vector obtained by averaging a minimum and a maximum
entry. In both cases $b$ is nonnegative, has the same total sum as $a$,
and satisfies
\[
 \sum_{|A|=k}\prod_{t\in A}a_t
 \le \sum_{|A|=k}\prod_{t\in A}b_t,
 \qquad
 \sum_t(b_t-c)^2\le\left(1-\frac1{2n}\right)\sum_t(a_t-c)^2.
\]
The subset sums range over subsets of $\{1,\ldots,n\}$.
\end{lemma}
\begin{proof}
If all entries are equal, take $b=a$. Then $c$ is that common entry,
both variances are zero, and both other claims are equalities. In
particular this treats $n=1$ without choosing distinct indices.
Otherwise choose indices $i,j$ attaining the minimum and maximum.
Their values satisfy $a_i<a_j$: otherwise the minimum equals the
maximum and every entry between them is equal, contradicting
nonconstancy. In particular $i\ne j$. Put $m=(a_i+a_j)/2$ and
$b=\operatorname{update}(\operatorname{update}(a,i,m),j,m)$.
The two entries at $i,j$ are $m$, and all other entries are unchanged,
which gives the operation certificate. Put $R=a_j-a_i>0$. By
\noderef{ElementarySymmetricPairAveraging}, $b$ is nonnegative, has the
same sum, its subset product sum does not decrease, and
\[
 \sum_t(b_t-c)^2=V-R^2/2,
 \qquad V=\sum_t(a_t-c)^2.
\]
Summing $a_i\le a_t\le a_j$ and dividing by $n>0$ gives
$a_i\le c\le a_j$. For every $t$ this implies
$-R\le a_t-c\le R$, hence $(a_t-c)^2\le R^2$.
Summing yields $V\le nR^2$. Since $2n>0$, we obtain
$V/(2n)\le R^2/2$. Consequently
\[
 \sum_t(b_t-c)^2=V-R^2/2
 \le V-V/(2n)=\left(1-\frac1{2n}\right)V,
\]
as claimed.
\end{proof}
